% !TEX root = paper.tex

\section{Related work}\label{sec:related}

Many tools for verifying Rust code exist.
As far as we know, no other tool leverages Rust's borrow checker to enforce linear ghost permissions.
However, in other dimensions, there is significant overlap between Verus and other projects.

Creusot~\cite{denis:hal-03737878} may be the closest tool to Verus,
since it uses Rust code to express specifications and proofs,
based on a macro named Pearlite.
Creusot functions can be annotated as \verb`#[logic]` or \verb`#[predicate]`
to indicate that the functions are ghost.
These are similar to Verus' \verb`spec` functions,
in that they are not checked for linearity and borrowing
(``Pearlite formulas are type-checked by the front-end of the Rust compiler,
but they are not borrow checked'').
Creusot does not have ghost code that is checked for linearity and borrowing,
the way Verus' \verb`proof` functions and \verb`proof` variables are.
Verus' SMT-LIB encoding is conceptually similar to the one produced 
by Creusot~\cite{denis:hal-03737878} via the Why3~\cite{boogie11why3} prover,
which requires an intermediate step: in Creusot the Rust code is first
lowered into Why3's MLCFG (an ML with labelled blocks and gotos),
and then Why3 encodes verification conditions for the backend solvers.

Prusti~\cite{DBLP:journals/pacmpl/Astrauskas0PS19,DBLP:conf/nfm/AstrauskasBFGMM22}
verifies Rust code by translating it into the Viper separation logic engine~\cite{DBLP:conf/vmcai/0001SS16},
effectively reverifying ownership properties enforced by Rust's borrow checker.
This relatively heavyweight encoding creates larger formulas for an SMT solver,
but can be used for Rust \verb`unsafe` code that subverts Rust's borrow checking rules.
By contrast, \verus relies on the memory safety enforced by Rust's borrow checker,
obviating the need to use separation logic ubiquitously---instead,
the user can selectively apply separation logic-style techniques
(based on linear ghost permissions) only for the tricky cases that require them.

Aeneas~\cite{DBLP:journals/pacmpl/HoP22} verifies Rust code by translating it into a purely functional
representation in \fstar.
In this style of verification,
programmers develop a proof about the functional representation of executable Rust code,
which is quite different from Verus' Hoare-logic style,
where the programmer annotates the Rust code with preconditions, postconditions, and loop invariants.

% {Verifying Rust Programs with SMACK} % https://soarlab.org/publications/2018_atva_bhr/

% {CRust: A Bounded Verifier for Rust} % https://alastairreid.github.io/RelatedWork/papers/toman:ase:2015/

\citet{DBLP:journals/pacmpl/YanovskiDJD21} propose a datatype called \verb`GhostCell`, which separates
data from permission in a manner similar to our \verb`PCell` and \verb`PermData`.
The main difference is that \verb`GhostCell` employs a polymorphic type trick to enforce that a permission may
only be used with the cells to which it corresponds, while \verb`PCell` uses a \verb`requires` clause to enforce this,
which is more flexible and allows permissions to depend on data that is not statically determined during type-checking.
Furthermore, while \verb`GhostCell` is used to enforce memory safety, to our knowledge, it has not been used to show functional correctness properties.

RustBelt~\cite{DBLP:journals/pacmpl/0002JKD18} is a verification framework that establishes
a semantic model for type safety in Rust: it allows a user to verify \verb`unsafe` code
with safe APIs, i.e., prove that any well-typed, \verb`unsafe`-free Rust program using the API
will be memory safe.
This makes it complementary to \verus, which relies on that memory safety, and indeed, it might be possible to use RustBelt to verify \verus' memory primitives (\verb`PPtr` and \verb`PCell`) and their specifications.
RustBelt can also handle atomics with relaxed memory ordering~\cite{DBLP:journals/pacmpl/DangJKD20}, which \verus does not support.
RustBelt is implemented in Coq, and thus proofs are written via tactics rather than by SMT.

RustBelt has also been used as part of RustHornBelt~\cite{DBLP:conf/pldi/MatsushitaDJD22}, which validates RustHorn~\cite{DBLP:conf/esop/0002T020},
the encoding used by Creusot.
However, RustHornBelt still requires that \verb`unsafe` code be proved correct in Coq,
while \verus provides safe, zero-cost alternatives to commonly used \verb`unsafe` Rust features via its linear ghost state.
Specifically, \verus provides \verb`PPtr` for raw pointers and \verb`PCell` for \verb`UnsafeCell`, so that users can write code (which would otherwise need those
\verb`unsafe` features) within \verus itself.

Note though that while \verus supports some \verb`unsafe` use-cases, including raw pointers,
our specification for pointers
is very simple, only handling pointers that point into heap allocations from the global
memory allocator. A complete pointer model for Rust would support pointers to the
stack variables, cell interiors, struct fields, references, and so on, as well as handle
thorny issues such as pointer provenance.
By comparison, Stacked Borrows~\cite{StackedBorrows} is a promising operational semantics 
for Rust memory accesses that aims to handle all these concepts.

Separation logic~\cite{DBLP:conf/lics/Reynolds02, DBLP:journals/tcs/OHearn07, Iris31}
was one inspiration for our linear ghost permissions,
although the techniques used in Verus and separation logic are quite different.
In separation logic, a permission is part of the logic rather than a program-level value,
and two permissions are combined together using separating conjunction.
In Verus, permissions are values and two permissions are combined together by placing them in a datatype.
Thus, in Verus, programmers manipulate permissions directly as data,
which can require extra programmer effort,
but makes generating verification conditions for an SMT solver much easier,
since SMT solvers handle classical logic, not separation logic.

Another inspiration for linear ghost permissions was earlier work on
using linearity in type systems to manage changing
state~\cite{crary1999capcalculus, smith2000aliastypes, zhu2005atsviews, morrisett2005l3}
Alias Types~\cite{smith2000aliastypes}, for example, tracks a set of constraints on the memory state,
and these constraints change linearly as the memory state evolves.
ATS~\cite{zhu2005atsviews} combines this idea, in the form of ``stateful views'',
with reasoning about integer arithmetic via a simple dependent type system.
Most similar to our approach is L3~\cite{morrisett2005l3},
which treats ``capabilities'' (permissions) as first-class linear ghost values,
as in Verus.  L3 uses type variables (specifically, location variables)
to connect the capabilities to pointers, whereas
Verus uses SMT solving, which avoids the burden on the programmer of
universally quantifying or existentially quantifying over location variables.
The combination of SMT solving and Rust's automated borrow checking means that
ideas from ATS and L3 are now not only possible within a mainstream language,
but convenient.

Dafny~\cite{DBLP:conf/lpar/Leino10} and \fstar~\cite{mumon} support ghost code and ghost variables.
\fstar uses an effect system to distinguish ghost functions from executable functions,
and has an \verb`erased` type to represent ghost data.
\fstar does not have a linear type system,
although the \fstar Steel system~\cite{DBLP:journals/pacmpl/FromherzRSGMMR21}
supports separation logic reasoning.
Dafny supports \verb`ghost` annotations on variables,
similar to Verus \verb`spec` variables,
and Dafny supports \verb`lemma`s, similar to Verus \verb`proof` functions.
Linear Dafny~\cite{DBLP:journals/pacmpl/LiLZCHPH22} extends Dafny with linear types and borrowing,
although the linearity and borrowing is less sophisticated than in Rust (for example,
Linear Dafny lacks lifetime variables).
